\section{Week 12: Gibbs Sampling and Convergence Diagnostics}

\begin{quote}
\textbf{Chapter 5: Parameter estimation with data-driven methods (MCMC)}
· Lecture 1 h · Lab 2 h \textbf{Running data:} log response times of one
participant (a normal model with unknown mean and precision), and the
power forgetting model sampled with last week's Metropolis sampler.
\end{quote}

\begin{center}\rule{0.5\linewidth}{0.5pt}\end{center}

\subsection{Key terms}

Important terms of this lesson, with their Persian equivalents (repeated
even if they appeared in earlier weeks).

\begin{longtable}[]{@{}
  >{\raggedright\arraybackslash}p{(\linewidth - 4\tabcolsep) * \real{0.3333}}
  >{\raggedright\arraybackslash}p{(\linewidth - 4\tabcolsep) * \real{0.3333}}
  >{\raggedright\arraybackslash}p{(\linewidth - 4\tabcolsep) * \real{0.3333}}@{}}
\caption{\centering Key terms}\tabularnewline
\toprule\noalign{}
\begin{minipage}[b]{\linewidth}\raggedright
English
\end{minipage} & \begin{minipage}[b]{\linewidth}\raggedright
\RL{\CMfa فارسی}
\end{minipage} & \begin{minipage}[b]{\linewidth}\raggedright
\RL{\CMfa توضیح کوتاه}
\end{minipage} \\
\midrule\noalign{}
\endfirsthead
\toprule\noalign{}
\begin{minipage}[b]{\linewidth}\raggedright
English
\end{minipage} & \begin{minipage}[b]{\linewidth}\raggedright
\RL{\CMfa فارسی}
\end{minipage} & \begin{minipage}[b]{\linewidth}\raggedright
\RL{\CMfa توضیح کوتاه}
\end{minipage} \\
\midrule\noalign{}
\endhead
\bottomrule\noalign{}
\endlastfoot
Gibbs sampling & \RL{\CMfa نمونه‌گیری} Gibbs & \RL{\CMfa به‌روزرسانی هر پارامتر از
توزیع شرطی کاملش} \\
convergence & \RL{\CMfa همگرایی} & \RL{\CMfa رسیدن زنجیره‌ها به توزیع هدف} \\
parameter & \RL{\CMfa پارامتر} & \RL{\CMfa کمیت قابل تنظیم مدل که بین افراد یا شرایط
فرق می‌کند} \\
estimate / estimator & \RL{\CMfa برآورد} / \RL{\CMfa برآوردگر} & \RL{\CMfa مقدار به‌دست‌آمده
برای پارامتر} / \RL{\CMfa روش به دست آوردن آن} \\
Markov chain Monte Carlo (MCMC) & \RL{\CMfa زنجیره‌ی مارکوف} Monte Carlo &
\RL{\CMfa نمونه‌گیری از پسین با یک زنجیره‌ی تصادفی} \\
response time (RT) & \RL{\CMfa زمان واکنش} & \RL{\CMfa فاصله‌ی ارائه‌ی محرک تا
پاسخ} \\
forgetting curve & \RL{\CMfa منحنی فراموشی} & \RL{\CMfa کاهش یادآوری با گذشت
زمان} \\
Metropolis algorithm & \RL{\CMfa الگوریتم} Metropolis & \RL{\CMfa پیشنهاد یک حرکت و
پذیرش یا رد آن} \\
sampler & \RL{\CMfa نمونه‌گیر} (sampler) & \RL{\CMfa الگوریتمی که از توزیع نمونه
می‌گیرد} \\
full conditional & \RL{\CMfa توزیع شرطی کامل} (full conditional) & \RL{\CMfa توزیع
یک پارامتر به شرط بقیه و داده} \\
conjugate prior & \RL{\CMfa پیشین مزدوج} (conjugate) & \RL{\CMfa پیشینی که پسین
هم‌خانواده می‌دهد} \\
R-hat & R-hat & \RL{\CMfa مقایسه‌ی واریانس بین و درون زنجیره‌ها؛ باید} \ensuremath{\approx} 1
\RL{\CMfa باشد} \\
effective sample size (ESS) & \RL{\CMfa اندازه‌ی مؤثر نمونه} (ESS) & \RL{\CMfa تعداد
نمونه‌ی مستقل هم‌ارز} \\
chain & \RL{\CMfa زنجیره} & \RL{\CMfa دنباله‌ی نمونه‌های} MCMC \\
trace plot & \RL{\CMfa نمودار} trace & \RL{\CMfa مقدار نمونه‌ها در برابر شماره‌ی
تکرار} \\
Monte Carlo & Monte Carlo \RL{\CMfa (مونت‌کارلو)} & \RL{\CMfa تقریب با میانگین
نمونه‌های تصادفی} \\
proposal & \RL{\CMfa پیشنهاد} (proposal) & \RL{\CMfa مقدار جدیدی که الگوریتم امتحان
می‌کند} \\
posterior & \RL{\CMfa توزیع پسین} (posterior) & \RL{\CMfa باور درباره‌ی پارامتر بعد
از دیدن داده} \\
MCSE & \RL{\CMfa خطای استاندارد} Monte Carlo & \RL{\CMfa خطای شبیه‌سازی در میانگین
پسین} \\
prior & \RL{\CMfa توزیع پیشین} (prior) & \RL{\CMfa باور درباره‌ی پارامتر پیش از دیدن
داده} \\
HMC / NUTS & Hamiltonian Monte Carlo / NUTS & MCMC \RL{\CMfa مبتنی بر
گرادیان} \\
burn-in / warm-up & \RL{\CMfa دوره‌ی گرم‌شدن} (burn-in) & \RL{\CMfa نمونه‌های اولیه که
دور ریخته می‌شوند} \\
autocorrelation & \RL{\CMfa خودهمبستگی} & \RL{\CMfa وابستگی نمونه‌های پشت سر هم} \\
sensitivity (d′) & \RL{\CMfa حساسیت} & \RL{\CMfa فاصله‌ی توزیع علامت و نویز} \\
\end{longtable}

\begin{center}\rule{0.5\linewidth}{0.5pt}\end{center}

\subsection{Learning outcomes}

By the end of this week, students will be able to:

\begin{enumerate}
\def\labelenumi{\arabic{enumi}.}
\tightlist
\item
  \textbf{Derive} the full conditional distributions of a semi-conjugate
  normal model and \textbf{formulate and evaluate} a Gibbs sampler from scratch.
\item
  \textbf{Explain} when Gibbs sampling is efficient and when it mixes
  slowly (correlated parameters), and how Metropolis-within-Gibbs
  handles non-conjugate conditionals.
\item
  \textbf{Diagnose} non-convergence from multiple chains using trace
  plots, \(\hat R\), ESS and the Monte Carlo standard error.
\item
  \textbf{Apply} a reporting checklist for MCMC results.
\end{enumerate}

\begin{center}\rule{0.5\linewidth}{0.5pt}\end{center}

\subsection{Motivation \& big picture}

Last week's Metropolis sampler had one tuning knob, the proposal SD, and
a nagging question: \emph{how do we know the chain has converged?} This
week answers both concerns.

\textbf{Gibbs sampling} removes the tuning problem for a large and
important class of models. Instead of proposing a jump for all
parameters and sometimes rejecting it, it updates \textbf{one parameter
at a time}, drawing it \emph{exactly} from its distribution given the
current values of all the others. Every draw is accepted. Many classic
Bayesian programs (BUGS, JAGS, used throughout Lee \& Wagenmakers and
Kruschke) are built on this idea.

\textbf{Convergence diagnostics} are the honest part of MCMC. A chain is
only useful if it has (a) forgotten its starting point and (b) explored
the whole posterior, and nothing in a single chain proves either. The
standard remedy is to run \textbf{several chains from different starting
points} and check that they agree. \(\hat R\) quantifies ``do the chains
agree?'', and the ESS quantifies ``how much independent information do
the correlated draws contain?''. No MCMC result should be reported
without them.

\paragraph{Where is this used?}

\begin{itemize}
\tightlist
\item
  \textbf{Every published MCMC analysis} must report convergence
  diagnostics; reviewers in psychology and neuroscience now routinely
  ask for \(\hat R\) and ESS.
\item
  \textbf{Gibbs sampling} powers the classic Bayesian programs BUGS and
  JAGS, which are used in many cognitive modeling textbooks (Lee \&
  Wagenmakers, Kruschke) and in ecology, epidemiology and education
  research.
\item
  \textbf{Image processing and communications}: Gibbs samplers have been
  used for image restoration and for detection in multi-antenna (MIMO)
  receivers.
\item
  \textbf{Topic models} (finding themes in large text collections) were
  classically fitted with Gibbs sampling.
\end{itemize}

\begin{center}\rule{0.5\linewidth}{0.5pt}\end{center}

\Needspace{11\baselineskip}

\subsection{Lecture notes}

\paragraph{In plain words}

Metropolis needs a well-chosen step size. \textbf{Gibbs sampling} avoids
that for many models by updating \textbf{one parameter at a time}: fix
all the other parameters at their current values, and draw the chosen
parameter directly from its exact distribution under those conditions
(its \textbf{full conditional}). Then move to the next parameter, and so
on. Nothing is ever rejected. The second half of the week answers the
uncomfortable question ``has my sampler actually worked?''. We run
\textbf{several chains} from very different starting points. If they all
end up describing the same distribution, that is good evidence of
convergence; \(\hat R\) (``R-hat'') is a single number that measures
whether the chains agree (it should be below 1.01). Because neighboring
draws are correlated, 4,000 draws may contain only the information of,
say, 1,000 independent ones: that number is the \textbf{effective sample
size (ESS)}, and it tells us how precise our averages are.

\textbf{Suggested timing (60 min):} Gibbs sampling 20 · when Gibbs
struggles 5 · convergence and multiple chains 8 · \(\hat R\) 10 · ESS
and MCSE 10 · reporting checklist 7.

\subsubsection{1. Gibbs sampling}

For parameters \(\boldsymbol\theta = (\theta_1, \dots, \theta_k)\), the
\textbf{full conditional} of \(\theta_m\) is its posterior given the
data and \emph{all other} parameters,
\(p(\theta_m \mid \boldsymbol\theta_{-m}, \mathbf{y})\). Gibbs sampling
cycles through them:

\begin{enumerate}
\def\labelenumi{\arabic{enumi}.}
\tightlist
\item
  Start at \(\boldsymbol\theta^{(0)}\).
\item
  For \(s = 1, \dots, S\):

  \begin{itemize}
  \tightlist
  \item
    draw
    \(\theta_1^{(s)} \sim p(\theta_1 \mid \theta_2^{(s-1)}, \dots, \theta_k^{(s-1)}, \mathbf{y})\)
  \item
    draw
    \(\theta_2^{(s)} \sim p(\theta_2 \mid \theta_1^{(s)}, \theta_3^{(s-1)}, \dots, \mathbf{y})\)
  \item
    \ldots{} (always using the newest values)
  \item
    draw
    \(\theta_k^{(s)} \sim p(\theta_k \mid \theta_1^{(s)}, \dots, \theta_{k-1}^{(s)}, \mathbf{y})\)
  \end{itemize}
\end{enumerate}

It is a special case of Metropolis--Hastings in which the proposal is
the full conditional itself, so the acceptance probability is always 1.

\textbf{Finding full conditionals.} Write the joint density
\(p(\mathbf{y}\mid\boldsymbol\theta)\,p(\boldsymbol\theta)\), keep only
the factors that contain \(\theta_m\), and recognize a known
distribution. This works when priors are (conditionally)
\textbf{conjugate}.

\textbf{Example: a normal model for log RTs.} Let \(y_i\) be log
response times, \(y_i \sim \mathcal{N}(\mu, 1/\sqrt{\tau})\), where
\(\tau = 1/\sigma^2\) is the \textbf{precision}. Priors:
\(\mu \sim \mathcal{N}(m_0, s_0)\) and
\(\tau \sim \mathrm{Gamma}(a_0, \text{rate } b_0)\), independent. Then:

\[
\mu \mid \tau, \mathbf{y} \sim \mathcal{N}\!\left(\frac{m_0/s_0^2 + \tau \sum_i y_i}{1/s_0^2 + N\tau},\ \ \frac{1}{\sqrt{1/s_0^2 + N\tau}}\right),
\]

\[
\tau \mid \mu, \mathbf{y} \sim \mathrm{Gamma}\!\left(a_0 + \frac{N}{2},\ \ \text{rate} = b_0 + \frac12\sum_i (y_i - \mu)^2\right).
\]

\textbf{Reading the formulas.} The conditional mean of \(\mu\) is a
precision-weighted average of the prior mean and the data mean (as in
the beta-binomial of Week 9). The conditional distribution of \(\tau\)
adds half an observation to the shape and half a squared residual to the
rate for each data point.

\begin{quote}
Mixing up rate and scale is the most common bug in Gibbs samplers.
\end{quote}

\paragraph{Reading the math}

\begin{itemize}
\tightlist
\item
  \(\boldsymbol\theta_{-m}\): ``all parameters except \(\theta_m\)''.
\item
  \textbf{Precision} \(\tau = 1/\sigma^2\): the inverse of the variance.
  Large precision = small noise. Working with precisions makes combining
  information a simple addition.
\item
  \textbf{The \(\mu\) update:} the conditional mean
  \(\dfrac{m_0/s_0^2 + \tau\sum y_i}{1/s_0^2 + N\tau}\) is a weighted
  average of the prior mean \(m_0\) (weight \(1/s_0^2\), the prior's
  precision) and the data mean (total weight \(N\tau\), the data's
  precision). The conditional precision is the \emph{sum} of the two
  precisions. Worked example: \(1 + 4\times10 = 41\), so the data
  dominate.
\item
  \textbf{The \(\tau\) update:} a Gamma distribution whose shape grows
  by \(1/2\) per observation and whose rate grows by half of each
  squared residual. Its mean, shape/rate, is roughly ``1 / (average
  squared residual)'', i.e.~the precision suggested by the data.
\end{itemize}

\subsubsection{2. When Gibbs struggles, and Metropolis-within-Gibbs}

\begin{itemize}
\tightlist
\item
  \textbf{Strong posterior correlations.} Gibbs moves parallel to the
  coordinate axes. If the posterior is a narrow diagonal ridge (like the
  \(A\)--\(d\) trade-off of the power model), each step can move only a
  tiny distance, and the chain zig-zags slowly along the ridge.
  Remedies: reparameterize, update correlated parameters jointly
  (``blocking''), or use gradient-based samplers (Week 13).
\item
  \textbf{Non-conjugate conditionals.} Most cognitive models (forgetting
  curves, SDT, the GCM) have no closed-form conditionals.
  \textbf{Metropolis-within-Gibbs} updates each parameter with a
  one-dimensional Metropolis step targeting its full conditional, and
  uses exact Gibbs draws where available.
\end{itemize}

\begin{quote}
\textbf{Advanced:} JAGS and BUGS automatically pick a sampler for each
parameter (conjugate Gibbs, slice sampling, adaptive Metropolis), which
is why they run on almost any model.
\end{quote}

\subsubsection{3. Convergence: the logic of multiple chains}

An MCMC chain has converged when its draws are representative of the
posterior. Two failure modes: 1. \textbf{Not yet stationary:} the chain
still reflects the starting value. 2. \textbf{Stuck or slow:} the chain
explores only part of the posterior (e.g., one of two modes, or one end
of a long ridge).

Neither can be detected reliably from one chain, because a chain stuck
in one region can look perfectly stable. Run \(C \ge 4\) chains from
\textbf{over-dispersed} starting points: if they all end up describing
the same distribution, that is good (though never conclusive) evidence
of convergence.

\textbf{Trace plots} of all chains overlaid should look like one
well-mixed ``hairy caterpillar'', with no chain visibly separated or
trending.

\subsubsection{\texorpdfstring{4. The potential scale reduction factor
\(\hat R\)}{4. The potential scale reduction factor \textbackslash hat R}}

Compare the variance \textbf{between} chains with the variance
\textbf{within} chains. With \(C\) chains of length \(S\) (after
burn-in), chain means \(\bar\theta_c\) and within-chain variances
\(s_c^2\):

\[
W = \frac{1}{C}\sum_c s_c^2, \qquad
B = \frac{S}{C-1}\sum_c (\bar\theta_c - \bar\theta)^2, \qquad
\widehat{\mathrm{var}}^{+} = \frac{S-1}{S}\,W + \frac{1}{S}\,B, \qquad
\hat R = \sqrt{\frac{\widehat{\mathrm{var}}^{+}}{W}}.
\]

If the chains have not mixed, \(B\) is large relative to \(W\) and
\(\hat R > 1\). As chains converge, \(\hat R \to 1\).

\begin{itemize}
\tightlist
\item
  \textbf{Split-\(\hat R\):} split each chain into a first and a second
  half and treat them as \(2C\) chains. This also detects a single chain
  that is still drifting.
\end{itemize}

\subsubsection{5. Effective sample size and Monte Carlo standard error}

Autocorrelated draws carry less information than independent ones. The
\textbf{effective sample size} is

\[
\mathrm{ESS} = \frac{C\cdot S}{1 + 2\sum_{h=1}^{\infty}\rho_h},
\]

where \(\rho_h\) is the lag-\(h\) autocorrelation, estimated from all
chains and truncated when the estimates become noise. ESS is the number
of \emph{independent} draws that would give the same precision.

\begin{itemize}
\tightlist
\item
  The \textbf{Monte Carlo standard error} of a posterior mean is
  \(\mathrm{MCSE} = \text{posterior SD}/\sqrt{\mathrm{ESS}}\). It must
  be small relative to the precision you report (e.g., don't report 3
  decimals if the MCSE is 0.01).
\item
  \textbf{Rule of thumb:} bulk-ESS \ensuremath{\geq} 400 (with 4 chains) for reliable
  posterior means and intervals; tail-ESS \ensuremath{\geq} 400 for reliable extreme
  quantiles.
\end{itemize}

\paragraph{Reading the math}

\begin{itemize}
\tightlist
\item
  \textbf{\(\hat R\):} \(W\) is the average variance \emph{inside} each
  chain; \(B\) measures how far the chain \emph{means} are from each
  other. \(\widehat{\mathrm{var}}^{+}\) is an estimate of the total
  posterior variance that includes the disagreement between chains. If
  the chains agree, \(B\) adds almost nothing and
  \(\hat R = \sqrt{\widehat{\mathrm{var}}^{+}/W} \approx 1\). In the
  ``not converged'' worked example, the two chains sit around 1.0 and
  2.0, \(B\) is huge and \(\hat R = 4.4\).
\item
  \textbf{ESS:} \(C \cdot S\) is the total number of draws; the
  denominator \(1 + 2\sum\rho_h\) inflates for correlated draws. With
  \(\rho_h = 0.5^h\), the sum is \(0.5 + 0.25 + \dots = 1\), so the
  draws count only one third: 4,000 draws are worth 1,333 independent
  ones.
\item
  \textbf{MCSE} \(= \text{posterior SD}/\sqrt{\mathrm{ESS}}\) is the
  simulation error of the reported mean; decide the number of decimals
  from it.
\end{itemize}

\subsubsection{6. MCMC reporting checklist}

\begin{enumerate}
\def\labelenumi{\arabic{enumi}.}
\tightlist
\item
  Sampler, number of chains, iterations, burn-in (warm-up), and any
  thinning.
\item
  Starting values (dispersed?).
\item
  \(\hat R\) (which version) and ESS for \textbf{all} reported
  parameters, plus the maximum \(\hat R\) and minimum ESS.
\item
  Trace plots (in a supplement) and a statement that they were
  inspected.
\item
  Priors, and the posterior summaries with MCSE-appropriate precision.
\item
  A posterior predictive check (Week 10).
\end{enumerate}

\paragraph{For an engineer}

\begin{itemize}
\tightlist
\item
  \textbf{Precision-weighted averaging} in the \(\mu\) update is the
  same rule as \textbf{maximal-ratio combining} or the Kalman
  measurement update: combine estimates in proportion to their inverse
  noise variances.
\item
  \textbf{Gibbs sampling} updates one variable given all the others,
  like coordinate-wise iterative algorithms (alternating optimization)
  and the iterative exchange of conditional information in turbo-style
  receivers; Gibbs samplers have been proposed for MIMO detection.
\item
  \textbf{\(\hat R\)} is a one-way ANOVA (between-group vs.~within-group
  variance) across independent simulation runs: the same check you would
  do by running a Monte Carlo simulation with several seeds and
  verifying that the runs agree.
\item
  \textbf{ESS} \(= N/(1 + 2\sum\rho_h)\) is the standard formula for the
  effective number of independent samples of a correlated random
  process: averaging \(N\) correlated samples reduces the variance only
  as much as averaging ESS independent ones (related to the correlation
  time).
\end{itemize}

\paragraph{For the instructor}

\begin{itemize}
\tightlist
\item
  \textbf{Draw Gibbs on a contour plot.} Two parameters, an elliptical
  posterior. Gibbs moves in horizontal and vertical steps only (a
  staircase). Then draw a narrow diagonal ellipse: the staircase steps
  become tiny. That one picture explains both how Gibbs works and when
  it struggles.
\item
  \textbf{Make \(\hat R\) visual first:} overlay four trace plots, once
  with chains that agree and once with one chain stuck elsewhere. Ask
  ``do these chains describe the same distribution?'' before showing any
  formula.
\item
  \textbf{Teach the reporting checklist as a habit.} Ask students to
  apply it to their project and to every MCMC result in the lab.
\item
  \textbf{Typical questions.}

  \begin{itemize}
  \tightlist
  \item
    \emph{``Should I thin my chains?''} Not for correctness; it only
    saves memory.
  \end{itemize}
\end{itemize}

\begin{center}\rule{0.5\linewidth}{0.5pt}\end{center}

\Needspace{11\baselineskip}

\subsection{Worked example}

\subsubsection{Part 1: one Gibbs iteration by hand}

Data: \(N = 4\) log RTs, \(\mathbf{y} = (6.0, 6.3, 6.1, 6.4)\), so
\(\sum y_i = 24.8\). Priors: \(m_0 = 6\), \(s_0 = 1\), \(a_0 = 2\),
\(b_0 = 0.1\). Current values: \(\mu = 6.0\), \(\tau = 10\).

\textbf{Update \(\mu\).} Precision \(= 1/1^2 + 4 \times 10 = 41\); mean
\(= (6/1 + 10 \times 24.8)/41 = (6 + 248)/41 = 6.195\); SD
\(= 1/\sqrt{41} = 0.156\). Draw, say, \(\mu = 6.22\).

\textbf{Update \(\tau\)} given \(\mu = 6.22\). Squared residuals:
\(0.0484 + 0.0064 + 0.0144 + 0.0324 = 0.1016\). Shape \(= 2 + 4/2 = 4\);
rate \(= 0.1 + 0.1016/2 = 0.1508\). The conditional mean of \(\tau\) is
\(4/0.1508 = 26.5\) (i.e., \(\sigma \approx 0.19\)). Draw, say,
\(\tau = 24.1\).

New state: \((\mu, \tau) = (6.22, 24.1)\). No accept/reject step was
needed.

\subsubsection{\texorpdfstring{Part 2: \(\hat R\) by
hand}{Part 2: \textbackslash hat R by hand}}

Two chains of \(S = 4\) draws each.

\textbf{Not converged.} Chain 1: (1.0, 1.2, 0.8, 1.0); chain 2: (2.0,
2.2, 1.8, 2.0). - Chain means: 1.0 and 2.0. Within-chain variances (ddof
= 1): both \(0.08/3 = 0.0267\), so \(W = 0.0267\). -
\(B = S \times \mathrm{var}(\text{means}) = 4 \times 0.5 = 2.0\). -
\(\widehat{\mathrm{var}}^{+} = \tfrac34 \times 0.0267 + \tfrac14 \times 2.0 = 0.520\).
- \(\hat R = \sqrt{0.520/0.0267} = 4.4\). The chains clearly disagree.

\textbf{Converged (realistic numbers).} \(S = 1000\) draws per chain,
\(W = 0.05\), chain means 1.00 and 1.02:
\(B = 1000 \times \mathrm{var}(1.00, 1.02) = 1000 \times 0.0002 = 0.2\);
\(\widehat{\mathrm{var}}^{+} = 0.999 \times 0.05 + 0.2/1000 = 0.05015\);
\(\hat R = \sqrt{0.05015/0.05} = 1.0015\). Below 1.01: fine.

\subsubsection{Part 3: ESS by hand}

\(C \cdot S = 4000\) draws with autocorrelations \(\rho_h = 0.5^h\).
Then \(\sum_{h\ge1}\rho_h = 1\), so
\(\mathrm{ESS} = 4000/(1 + 2) = 1333\). With a posterior SD of 0.2, MCSE
\(= 0.2/\sqrt{1333} = 0.0055\): report the posterior mean to two
decimals.

\begin{center}\rule{0.5\linewidth}{0.5pt}\end{center}

\Needspace{11\baselineskip}

\subsection{Lab (2 hours)}

\subsubsection{Goals}

\begin{enumerate}
\def\labelenumi{\arabic{enumi}.}
\tightlist
\item
  Formulate and evaluate a Gibbs sampler for the normal model of log RTs and validate
  it against a grid posterior.
\item
  Run four dispersed chains and inspect them.
\item
  Break convergence on purpose (the power model with a bad Metropolis
  sampler) and watch the diagnostics catch it.
\end{enumerate}

\textbf{Timing:} Part A 30 min · Part B 15 min · Part C 30 min · Part D
30 min · Part E 15 min.

\subsubsection{Part B: four dispersed chains (15 min)}

Gibbs forgets even absurd starting values within a few iterations here,
because \(\mu\) and \(\tau\) are nearly uncorrelated in this posterior.

\subsubsection{Part C: R-hat and ESS diagnostics (30 min)}

With 4 × 2500 = 10 000 draws, ESS \ensuremath{\approx} 9000--10 000 means the Gibbs draws
are almost independent.

\subsubsection{Part D: catching non-convergence (30 min)}

We reuse last week's Metropolis sampler on the power forgetting model,
once badly tuned and short, once well tuned.

\textbf{What to notice:} with tiny steps and short chains, each chain is
still close to where it started. The traces are separate, \(\hat R\) is
far above 1.01 and the ESS is tiny. \emph{Each chain on its own looks
smooth and stable}, which is exactly why one chain is never enough. The
tuned sampler gives overlapping traces, \(\hat R = 1.002\) and an ESS of
about 1000 from 20 000 draws. The Gibbs sampler of Part C, with almost
uncorrelated draws, reached an ESS of about 9000--10 000 from 10 000
draws.

\subsubsection{Part E: interpret and report (15 min)}

Write an MCMC methods-and-results paragraph for the log-RT model,
following the checklist in Lecture §6: sampler, chains, iterations,
burn-in, starting values, \(\hat R\), ESS, posterior means and 95\%
intervals of \(\mu\) and \(\sigma\) (back-transform \(\mu\) to
milliseconds with \(e^\mu\) and say what it represents).

\subsubsection{Exercises (increasing difficulty)}

\textbf{Exercise 1 \ensuremath{\star}.} Show by algebra that the full conditional of
\(\tau\) is a Gamma distribution: write
\(p(\tau\mid\mu,\mathbf{y}) \propto \tau^{N/2}\exp(-\tfrac{\tau}{2}\sum(y_i-\mu)^2)\cdot\tau^{a_0-1}e^{-b_0\tau}\)
and read off the shape and rate.

Do the diagnostics catch it? Does the validation against the grid? What
does this tell you about the difference between \emph{convergence} and
\emph{correctness}?

\textbf{Exercise 3 \ensuremath{\star}\ensuremath{\star} (correlated Gibbs).} Write a Gibbs sampler for a
bivariate normal ``posterior'' with correlation \(\rho\) (full
conditionals:
\(\theta_1\mid\theta_2 \sim \mathcal{N}(\rho\theta_2, \sqrt{1-\rho^2})\)
and vice versa). Compare the ESS for \(\rho = 0, 0.9, 0.99\) with 5000
iterations, and plot the zig-zag path for \(\rho = 0.99\).

\textbf{Exercise 4 \ensuremath{\star}\ensuremath{\star} (Metropolis-within-Gibbs).} Sample the power
forgetting model by updating \(A\) and \(d\) \textbf{one at a time},
each with a 1-D Metropolis step. Compare the ESS per second with the
joint Metropolis sampler. Why might it be worse here? (Hint: the
\(A\)--\(d\) correlation.)

\textbf{Exercise 5 \ensuremath{\star}\ensuremath{\star}\ensuremath{\star} (Advanced: two conditions).} Extend the Gibbs
sampler to log RTs from two conditions with separate means
\(\mu_1, \mu_2\) and a shared precision \(\tau\). Simulate data with a
true difference of 0.1, sample, and report the posterior of
\(\mu_2 - \mu_1\) and \(P(\mu_2 > \mu_1\mid\mathbf{y})\) with
diagnostics.

\begin{center}\rule{0.5\linewidth}{0.5pt}\end{center}

\subsection{Connection to AI \& LLMs}

\textbf{Gibbs sampling from a masked language model.} Models like BERT
are trained to fill in a masked word given all other words: they learn
\textbf{full conditionals}, \(P(w_i \mid \text{all other words})\). That
is exactly what a Gibbs sampler needs. Repeatedly masking one position
and resampling it generates whole sentences (Wang \& Cho, 2019). The toy
model below knows only the two full conditionals of a two-word
``language'', and Gibbs sampling recovers the joint distribution of
sentences. Because the language is ``sticky'' (dog goes with barks, cat
with meows), the chain mixes slowly, and \(\hat R\) catches this in
short runs.

How many iterations do you now need for \(\hat R < 1.01\)? - \ensuremath{\star}\ensuremath{\star} Real
masked models are trained separately for each position, so their
conditionals need not come from any single joint distribution. What does
this tell you about trusting a sampler whose conditionals are learned?

\emph{Optional reading:} Wang, A., \& Cho, K. (2019). BERT has a mouth,
and it must speak: BERT as a Markov random field language model.
\emph{arXiv:1902.04094}.

\begin{center}\rule{0.5\linewidth}{0.5pt}\end{center}

\subsection{Common mistakes \& misconceptions}

\begin{enumerate}
\def\labelenumi{\arabic{enumi}.}
\tightlist
\item
  \textbf{Using old values within a Gibbs sweep.} Each conditional must
  use the \emph{newest} values of the other parameters.
\item
  \textbf{One long chain instead of several.} A stuck chain can look
  perfectly stable. Use \(\ge 4\) dispersed chains.
\item
  \textbf{Reporting results with low ESS.} With ESS = 50, a posterior
  mean has an MCSE of about 14\% of the posterior SD, and tail quantiles
  are unreliable.
\item
  \textbf{Thinning to ``fix'' autocorrelation.} Thinning lowers storage
  but does not increase the information. ESS already accounts for
  autocorrelation.
\item
  \textbf{Using the 1.1 threshold uncritically.} Modern guidance is
  \(\hat R < 1.01\) with rank-normalized split-\(\hat R\).
\item
  \textbf{Diagnosing only some parameters.} Check \emph{all} reported
  parameters, including derived ones.
\end{enumerate}

\begin{center}\rule{0.5\linewidth}{0.5pt}\end{center}

\subsection{Check your understanding}

\textbf{1.} Why does Gibbs sampling never reject a draw? \textbf{Answer}
Each draw comes exactly from the full conditional distribution. Viewed
as Metropolis--Hastings with the full conditional as proposal, the
acceptance ratio is always 1.

\textbf{2.} What do you need to be able to run a pure Gibbs sampler?
\textbf{Answer} Full conditional distributions for every parameter (or
block) that you can sample from directly, typically because the priors
are conditionally conjugate.

\textbf{3.} A Gibbs sampler for two strongly correlated parameters has a
very low ESS. Why, and what could you do? \textbf{Answer} Gibbs updates
one coordinate at a time, so on a narrow diagonal ridge every step is
tiny and the chain zig-zags slowly. Remedies: reparameterize to reduce
the correlation, update the correlated parameters jointly (blocking), or
use HMC/NUTS.

\textbf{4.} Explain \(\hat R\) in one sentence. Which value indicates
convergence? \textbf{Answer} It compares the variance between chains
with the variance within chains (it estimates how much the posterior
spread could shrink if the chains ran forever); values near 1
(\textless{} 1.01 by current guidance) indicate that the chains agree.

\textbf{5.} You have 4000 draws and an ESS of 200 for a parameter with
posterior SD 0.5. What is the MCSE of its posterior mean? How many
decimals should you report? \textbf{Answer} MCSE
\(= 0.5/\sqrt{200} = 0.035\). Report the mean to one decimal (or two at
most), because the third decimal is Monte Carlo noise.

\textbf{6.} Why can't one chain demonstrate convergence? \textbf{Answer}
A chain stuck in one region (one mode, or one part of a ridge) can look
stationary and well mixed, and nothing in the chain itself reveals the
unexplored regions. Only chains started elsewhere, which end up
somewhere different, expose the problem.

\textbf{7.} Name four things a methods section should report about an
MCMC analysis.

\begin{center}\rule{0.5\linewidth}{0.5pt}\end{center}

\Needspace{11\baselineskip}

\subsection{Assignment}

\subsubsection{A5: MCMC from scratch, with diagnostics (1.0 point of 20;
released Week 12, due at the start of the Week 14 lecture)}

\textbf{Tasks.}

\begin{enumerate}
\def\labelenumi{\arabic{enumi}.}
\tightlist
\item
  \textbf{Metropolis for SDT (3 pts).} Formulate and evaluate a random-walk
  Metropolis sampler for the shared-\(d'\) payoff model (\(d'\),
  \(c_1, c_2, c_3\)) with priors \(d' \sim \mathcal{N}(1, 1.5)\) and
  \(c_m \sim \mathcal{N}(0, 1)\). Tune the proposal (report the
  acceptance rate), and run 4 dispersed chains.
\item
  \textbf{Inference (1.5 pts).} Report posterior means and 95\%
  intervals; compute the posterior of
  \(c_{\text{conservative}} - c_{\text{liberal}}\) and the probability
  that it is positive. Compare \(d'\) with the Week 4 MLE.
\item
  \textbf{Gibbs (2 pts).} Formulate and evaluate the two-condition Gibbs sampler of
  Exercise 5 and report the posterior of the mean difference with
  diagnostics.
\item
  \textbf{Interpretation (1.5 pts, max. 300 words).} What do the payoffs
  do to sensitivity and bias, according to the posterior? Compare with
  the likelihood-ratio conclusion of Week 4.
\end{enumerate}

\textbf{Rubric:} points as above; within each task: correctness 60\%,
figures/tables 20\%, interpretation 20\%. Late submissions lose 10\% per
day (max. 3 days).

\begin{center}\rule{0.5\linewidth}{0.5pt}\end{center}

\subsection{Further reading}

\begin{itemize}
\tightlist
\item
  \textbf{Kruschke (2014), chapter 7} (sections on Gibbs sampling and on
  MCMC representativeness, accuracy and efficiency): diagnostics, ESS
  and MCSE explained visually. \emph{(Main reading.)}
\item
  \textbf{Lee \& Wagenmakers (2014), chapter 4.} Inferences with
  Gaussians: the normal model whose Gibbs sampler you built today, in
  its Bayesian-modeling context.
\item
  \emph{Optional, outside the course list:} Vehtari, A., Gelman, A.,
  Simpson, D., Carpenter, B., \& Bürkner, P.-C. (2021).
  Rank-normalization, folding, and localization: An improved \(\hat R\)
  for assessing convergence of MCMC. \emph{Bayesian Analysis, 16}(2),
  667--718.
\end{itemize}

\emph{Preview for Week 13:} writing samplers by hand teaches how they
work, but it is slow and error-prone.
